\newpage
\section{操作系统}
在本项目的设计和建设过程中，针对不同的应用场景，
引入多种不同的操作系统。

\subsection{桌面操作系统}
桌面操作系统为研究者使用的终端设备，
因此对环境的要求不大，满足日常使用即可。

\subsubsection{Windows\cite{windows11}}
Windows 为微软公司开发的操作系统，安装于大多数个人及企业
所购置的计算机设备上，兼容广泛。
同时也是大部分商业软件的目标平台之一。

\subsubsection{MacOS\cite{macOSSequoia}}
MacOS 为苹果公司开发的操作系统，主要安装于苹果公司推出的设备上，以良好的生态著称。

\subsubsection{linux发行版}
基于linux内核的操作系统发行版：

\paragraph{ubuntu\cite{ubuntuChina}}
Canonical公司推出的开源操作系统，有极强的兼容性（附带大多数驱动），
以用户友好著称。

\paragraph{fedora\cite{fedoraProject}}
RedHat公司推出的开源操作系统。

\paragraph{debian\cite{debian}}
部分桌面环境使用此操作系统。
因为软件包更新较慢，所以大多运行与相对较旧的设备上

\paragraph{arch linux\cite{archLinux}}
部分桌面环境使用此操作系统。此操作系统使用结构简单的包管理器。
大部分的复杂操作由用户完成。使用systemd，默认不启用任何服务，全部由用户决定。自定义程度较高。

\paragraph{gentoo\cite{gentoo}}
部分桌面环境使用此操作系统。此操作系统推荐从源代码构建所有软件包。
有一套极为完善的源代码构建系统，自定义程度极高。

\subsubsection{*bsd}
MacOS的核心版本Darwin基于freebsd等类BSD操作系统，
但是基于BSD系统配置开源图形化界面即为困难，
外加上大部分开源软件的生态都是基于Linux系列操作系统的，
本项目未涉及BSD操作系统。

\subsection{服务器操作系统}

\subsubsection{Windows}
Windows 有特定的服务器版本。内置服务器需要的各种组件。

\subsubsection{MacOS}
苹果公司与2022年停止了MacOS server的开发，
其他相关支持已于October 29, 2024,停止。

\subsubsection{linux发行版}
linux 发行版本大多数兼容POSIX标准，因此非常适合作为服务器使用。

\paragraph{ubuntu}
ubuntu的用户友好在服务器端也有良好的体现，其中预置的snap更带来了开箱即用的体验。

\paragraph{fedora}
红帽系列的发行版。

\paragraph{debian}
因为软件包更新周期很长，所以维护频率也不高。

\paragraph{arch linux}
archlinux更新非常频繁，所以不是很适合作为服务器系统，
通常需要引入冗余备份才能够稳定使用。并且需要切换至更新不快的lts内核，
否则需要经常重启，周期最短可至一天一次。

\paragraph{gentoo}
gentoo可以选择稳定分支的代码仓库，并且可以根据使用特定版本的软件包。
虽然需要很长的时间上线，但是足够稳定。
在相同型号设备很多的时候可以考虑。

\subsection{嵌入式操作系统}

\subsubsection{Windows}
微软针对嵌入式设备提供windows IoT操作系统。

\subsubsection{ubuntu}
大多数嵌入式设备都会优先针对类debian操作系统适配操作系统。
因此ubuntu通常是具备硬件厂商官方支持的发行版。

\subsubsection{debian}
同上